\chapter{Protocollo Sperimentale}

\section{Istanze}

Gli esperimenti sono condotti su 10 istanze dal benchmark CVRPLIB \cite{cvrplib}, suddivise in quattro
set con caratteristiche differenti:

\begin{table}[H]
\centering
\caption{Istanze CVRPLIB utilizzate}
\label{tab:instances}
\begin{tabular}{@{}llrrr@{}}
\toprule
\textbf{Set} & \textbf{Istanza} & \textbf{Nodi} & \textbf{Veicoli} & \textbf{Ottimo} \\
\midrule
A & A-n45-k7  & 45  &  7 & 1{,}146 \\
A & A-n60-k9  & 60  &  9 & 1{,}354 \\
A & A-n80-k10 & 80  & 10 & 1{,}763 \\
\midrule
B & B-n56-k7  & 56  &  7 &   707 \\
B & B-n66-k9  & 66  &  9 & 1{,}316 \\
B & B-n78-k10 & 78  & 10 & 1{,}221 \\
\midrule
E & E-n76-k8  & 76  &  8 &   735 \\
E & E-n101-k14& 101 & 14 & 1{,}071 \\
\midrule
P & P-n50-k10 & 50  & 10 &   696 \\
P & P-n101-k4 & 101 &  4 &   681 \\
\bottomrule
\end{tabular}
\end{table}

I quattro set presentano caratteristiche distintive:
\begin{itemize}
    \item \textbf{Set A}: clienti distribuiti uniformemente, domande omogenee;
    \item \textbf{Set B}: clienti concentrati in cluster geografici;
    \item \textbf{Set E}: distribuzione mista con cluster asimmetrici ---
          considerato il set pi\`u difficile;
    \item \textbf{Set P}: clienti con domanda molto variabile (fino a 4 veicoli
          per 101 clienti).
\end{itemize}

\section{Setup Sperimentale}

Per ogni istanza vengono eseguiti $R = 5$ run indipendenti con semi casuali diversi.
Per ogni run vengono misurati:
\begin{itemize}
    \item Costo della migliore soluzione trovata;
    \item Numero di generazioni per raggiungere la migliore soluzione;
    \item Storico di convergenza (best cost nel tempo, campionato ogni 100 FE);
    \item Numero di veicoli utilizzati e route complete.
\end{itemize}

Vengono poi calcolate le statistiche aggregate: best, mean, standard deviation e gap
percentuale dal Best Known Solution (BKS). Il budget computazionale \`e fissato a
$FE = 350{,}000$ valutazioni per ogni run.

\section{Configurazione Tuned}

La configurazione principale \texttt{config\_tuned} \`e stata ottenuta ottimizzando
i parametri dell'HGA tramite Optuna \cite{akiba2019optuna}, un framework di ottimizzazione iperparametrica
bayesiana. L'ottimizzazione ha esplorato lo spazio dei parametri su 100 trial,
valutando ogni configurazione su un subset di 4 istanze rappresentative (una per set)
con 3 run ciascuna e un budget ridotto di 100{,}000 FE per trial.

I parametri ottimizzati sono: dimensione della popolazione $\mu \in [5, 100]$,
tasso di crossover $p_c \in [0.5, 1.0]$, tasso di mutazione $p_m \in [0.05, 0.5]$,
tasso di ricerca locale $p_{ls} \in [0.05, 0.5]$, tournament size $k \in [2, 6]$,
elite count $e \in [1, 10]$ e granular size $\gamma \in [2, 40]$.

Il risultato dell'ottimizzazione ha prodotto una configurazione con $\mu = 81$,
$p_c = 0.675$, $p_m = 0.236$, $p_{ls} = 0.259$, $k = 4$, $e = 4$, $\gamma = 25$,
che bilancia efficacemente esplorazione globale e raffinamento locale.
